\answer{ex:19:1}{(а)~Один синапс даёт $6.7\mV$ ($R=66.7$~МОм), так что до порога вообще нужно не меньше $4$ (трёх хватает лишь асимптотически); $8$ за $10\ms$; $14$ за $5\ms$. (б)~$0.1\ms\ll\tau_m$, поправка от утечки $\sim0.25\%$.}
\answer{ex:19:2}{(а)~$1.75\cdot10^{10}$ (четверть смесителей отвечает на всё подряд), $4.4\cdot10^9$ и $0.06$ (при $k=40$ не отвечает почти никогда). (б)~В $10^3$ раз: $2\cdot10^5$ против $200$.}
\answer{ex:19:3}{(а)~$C(2n,n)=2^{2n-1}$ по симметрии биномиальных коэффициентов. (б)~$0.93$ и $0.09$; ширина перехода $\sim\sqrt{2n}$ по $P$, относительная $\sim1/\sqrt n$.}
\answer{ex:19:4}{С одного дендрита $V_{\text{тело}}<\kappa E_E<\theta$ при любом числе входов; при $g_0=0.5g_L$ нужно $n\ge3$ на каждом дендрите.}
\answer{ex:19:5}{$I\ge C\sigma_V/\sigma_t=15$~нА, то есть $150$ синхронных синапсов --- нереалистично; маленькому нейрону хватило бы $1$~нА, а с синапсом в $4$~нА дрожание $5$~мкс.}
\answer{ex:19:6}{Максимум на ближнем конце: $0.03\,\tau_m$ и $0.97$ ($L=0.2$); $0.32\,\tau_m$ и $0.66$ ($L=1$); $0.79\,\tau_m$ и $0.33$ ($L=2$). Оценка~\eqref{eq:dendriteload} по полной ёмкости верна только при $L\ll1$.}
\answer{ex:19:7}{Открытая задача. При фиксированном $m$ время отклика растёт линейно с запоминаемым $P$; при фиксированной доле $m=\varphi n$ оно перестаёт зависеть от $n$, и медленность большого нейрона --- следствие суммирования многих входов, а не размера как такового.}
